\begin{claim}{F.7}{Open}{The nonsolvable centralizer comparison and its checked finite image}
For an oriented mapping class group $M$ of genus at most two,
the normalizer and centralizer comparisons for finite solvable
subgroups have the complete proof below. For a geometric
$A_5\le M=\Gamma^+_{0,\nu}$, the equality
\[
                 C_{\widehat M}(A_5)
                    =\overline{C_M(A_5)}                 \tag{F.7.1}
\]
is stated in \cite[Theorem~C(iii)]{R-Boggi} and
\cite[Theorem~5.1(iii)]{R-BZ}. The displayed proofs contain
an additional closure--intersection equality for which we
have not verified a justification. We do not assert that
(F.7.1) is false.

There is nevertheless equality of their images in the
finite puncture-permutation quotient. Write the invariant
puncture set as a union of some of the exceptional orbits
of sizes $12,20,30$ and $m$ free orbits of size $60$.
Both images fix the exceptional punctures individually
and induce the full group $A_5^m\rtimes S_m$ on the
free orbits. In particular, with only exceptional orbits
the unresolved part of (F.7.1) lies in the pure profinite
mapping class group; its discrete right side is trivial.
\end{claim}
\begin{proof}[Source audit and proved comparisons]
The supplied preprint snapshot of \cite{R-Boggi}
has since appeared in \emph{S\~ao Paulo Journal of
Mathematical Sciences}, volume~20, article~35,
published 23 September 2026, DOI
\texttt{10.1007/s40863-026-00558-y}.
The publisher's full text was checked as well:
Proposition~2.4 retains the same displayed
closure--intersection step, and Lemma~3.2
still invokes that proposition.
The audit below applies to both inspected versions;
it is a proof-step audit and makes no claim based
on the paper's publication status.

For every finite subgroup $J\le M$, there are canonical
continuous maps
\[
 \widehat{C_M(J)}\longrightarrow
       \overline{C_M(J)}\lhook\joinrel\longrightarrow
       C_{\widehat M}(J).                              \tag{F.7.2}
\]
The first is onto by density and compactness. Its
injectivity states that the topology induced from $M$
on the discrete centralizer is its full profinite
topology. Surjectivity of the second arrow is the
ambient centralizer comparison. These are two specific
properties of the displayed maps; injectivity of the
first arrow does not supply surjectivity of the second.

In \cite[Proposition~2.4]{R-Boggi}, choose the displayed
generators $x_1,\ldots,x_k$ of the finite group, each of
prime-power order. The prime-power centralizer theorem
does prove
\[
 C_{\widehat M}(J)
   =\bigcap_i C_{\widehat M}(x_i)
   =\bigcap_i\overline{C_M(x_i)}.                      \tag{F.7.3}
\]
The source then identifies this last intersection with
$\overline{\bigcap_i C_M(x_i)}$. The same final equality
occurs in the proof of
\cite[Theorem~5.1(iii)]{R-BZ}, there indexed by all
$x\in J$. Neither displayed argument supplies a
closure--intersection hypothesis. This exact step is
what is left unverified here. The later virtually
compact special application in \cite[\S6.2]{R-BZ}
has a separate relatively quasiconvex product-separability
input; no such input has been proved here for the
punctured-sphere mapping class group.

\emph{A noncircular solvable proof.}
We first record the prime-power inputs used in the
induction. Goodness and finite virtual cohomological
type of mapping class groups of genus at most two,
and their finite-index subgroups, permit the direct
application of \cite[Theorem~A]{R-BZ}. Thus for a
finite $p$-subgroup $A$ we have
\[
 N_{\widehat M}(A)=\overline{N_M(A)},\qquad
 C_{\widehat M}(A)=\overline{C_M(A)}.                   \tag{F.7.4}
\]
The proof of \cite[Theorem~A and Corollary~B]{R-Boggi}
also supplies the full induced profinite topology on
$C_M(A)$ and hence on $N_M(A)$ in this prime-power case.
Here is why this use avoids the preceding gap. Its
free-action argument is the restriction-injectivity
argument on the finite covering surface group. Its
ramification-removal argument uses Proposition~2.4
only for the specified group $A$, now a $p$-group,
in the good configuration-group extension. For that
application the second equality in (F.7.4), directly
from \cite[Theorem~A]{R-BZ}, replaces Proposition~2.4's
intersection argument. Thus the two required inputs
are justified without using any nonsolvable assertion.
The passage from centralizer to normalizer topology
uses its finite-index kernel under conjugation into
$\Aut(A)$; this kernel is open in the induced topology,
as is detected by any finite quotient in which $A$
injects.

We prove the solvable normalizer comparison, together
with its full induced topology, by induction on the
order of the finite solvable group, simultaneously for
all finite-index subgroups of these mapping class
groups. Finite-index passage preserves the comparison:
intersect the completed normalizer with the corresponding
open subgroup and use density in an open set.

For a nontrivial finite solvable $J$, take a nontrivial
$p$-primary subgroup $A$ of the last nonzero term of
its derived series. That term is abelian, so its
$p$-primary subgroup is characteristic in $J$.
Consequently every normalizer of $J$ normalizes $A$.
The discrete covering normalizer sequence is
\[
 1\longrightarrow A\longrightarrow N_M(A)
                 \longrightarrow Q\longrightarrow1,   \tag{F.7.5}
\]
where $Q$ is a finite-index subgroup of the mapping
class group of the quotient surface with its branch
and puncture decorations. Its genus is again at most
two. The exact quotient map is induced by descent of
the commuting covering diagram; the decorations
record the original punctures, the removed branch
points and the covering subgroup. This is the discrete
normalizer sequence in \cite[\S2.1]{R-Boggi}.

By the prime-power topology input, the completed
normalizer in (F.7.4) is $\widehat{N_M(A)}$.
Completing (F.7.5) therefore gives its quotient by
the finite embedded kernel $A$ as $\widehat Q$.
The subgroup $J/A\le Q$ is smaller and solvable.
The induction hypothesis applies to its normalizer
in $\widehat Q$. The inverse image of that normalizer
under the quotient in (F.7.5), on either side, is
exactly the normalizer of $J$: $J$ is the full inverse
image of $J/A$ because it contains $A$. Pulling back
the dense discrete normalizer proves
\[
                  N_{\widehat M}(J)=\overline{N_M(J)}.
                                                               \tag{F.7.6}
\]
It also proves the topology assertion. Indeed the
completion of this inverse image maps to
$\widehat{N_M(A)}$; on the quotient by $A$ it is
injective by induction, and on $A$ it is the specified
inclusion. Its kernel is therefore trivial.
This completes the induction.

Conjugation gives a continuous homomorphism
\[
       N_{\widehat M}(J)\longrightarrow\Aut(J)
\]
to a finite group. Its kernel is the ambient centralizer
and is open. The dense subgroup $N_M(J)$ from (F.7.6)
meets that open kernel densely, and its intersection
there is exactly $C_M(J)$. This proves
\[
                  C_{\widehat M}(J)=\overline{C_M(J)}
\]
for every finite solvable $J$, without taking an
unjustified intersection of closures.

\emph{The finite image for $A_5$.}
Realize the geometric $A_5$ as the icosahedral rotation
group. Its exceptional point stabilizers have orders
$2,3,5$, and the corresponding transitive orbits are
$A_5/C_2$, $A_5/C_3$, $A_5/C_5$. Every permutation
commuting with this transitive action is right
multiplication by an element of $N_{A_5}(C_d)/C_d$.
This quotient has order two in each of the three
cases; its nonidentity permutation exchanges the
two poles of each relevant rotation axis.

An orientation-preserving homeomorphism commuting
with an odd-order nonidentity rotation fixes its
two poles individually. The local oriented rotation
numbers at the two poles are opposite, and an
orientation-preserving commuting homeomorphism
preserves this local rotation number. This excludes
the nonidentity permutation on the $12$- and
$20$-point orbits for the discrete cyclic centralizers.
Their completed centralizers have the same finite
permutation images by (F.7.4). The full $A_5$
centralizer lies in these cyclic centralizers, so
it fixes these exceptional punctures individually.

For an order-two stabilizer choose a tetrahedral
subgroup $A_4\le A_5$ containing it. An
orientation-preserving homeomorphism commuting
with $A_4$ fixes every ramification point of the
$A_4$ quotient. To see this, its quotient fixes the
branch point of order two. It also fixes each of the
two branch points of order three, rather than
exchanging them, by the preceding local rotation
argument. Isotope the quotient map to the identity
relative to these three branch points and lift the
isotopy. The lift stays in the centralizer; its
endpoint is a deck transformation in $Z(A_4)=1$.
The ramification set is finite and invariant along
the isotopy, so every one of its points is fixed.
The usual discrete covering-centralizer comparison
provides a commuting representative for a centralizing
mapping class. Since $A_4$ is solvable, the proved
centralizer comparison transfers this fixed-puncture
property to $C_{\widehat M}(A_4)$. It therefore holds
for $C_{\widehat M}(A_5)$, excluding the last possible
nonidentity permutation on the $30$-point orbit.

On $m$ free orbits the finite commutant is
$A_5^m\rtimes S_m$. Every element of this commutant
is already realized by the discrete centralizer.
For a direct construction, remove the three branch
points from the quotient sphere. Its unramified
$A_5$-cover has surjective monodromy onto $A_5$.
Move one of the $m$ quotient punctures around a
loop, keeping all the others fixed, and lift the
isotopy starting at the identity. The lift commutes
with $A_5$ and acts on that one free orbit by its
specified monodromy, while fixing the other orbits
and all ramification points. These loops realize
each factor $A_5$. Half-twists interchanging the
quotient punctures realize $S_m$ modulo these factors.
Thus the displayed wreath product is the exact
finite permutation image on both sides of (F.7.1).

The program
\path{code/part_a/a5_centralizer_permutations.py}
constructs $A_5$ as the sixty even permutations of
five letters. It checks its order-two and order-three
generating pair, the three coset actions, normalizer
orders $4,6,10$, and commutants of order two. The
fixed-pole constraints leave only the identity in
each exceptional block; the order-two block uses
the six ramification points for the specified $A_4$.

For completeness, the residual intersection question
has an exact finite-quotient formulation. Take
$u=(1\ 2)(3\ 4)$ and $v=(1\ 3\ 5)$ in $A_5$, which
generate it, and put $D_u=C_M(u)$, $D_v=C_M(v)$.
For open normal subgroups $V\subset U$ of $M$, let
$D_u(V),D_v(V)$ denote their images in $M/V$.
Compactness of coset fibres proves
\[
 \operatorname{im}\bigl(C_{\widehat M}(A_5)\to M/U\bigr)
 =\bigcap_{V\subset U}
       \operatorname{im}\bigl(D_u(V)\cap D_v(V)\to M/U\bigr).
                                                               \tag{F.7.7}
\]
Indeed choose a lift $t\in\widehat M$ of a point
in the right side. For each $V\subset U$, the set
\[
 (\overline{D_u}\,\overline V)\cap
 (\overline{D_v}\,\overline V)\cap t\overline U
\]
is nonempty and compact. These sets are directed
under refinement, so their finite intersections
are nonempty. Compactness produces a single element
in their intersection, which is
$\overline{D_u}\cap\overline{D_v}\cap t\overline U$.
Conversely any such element supplies every finite
coordinate. Equations (F.7.3) and (F.7.7) identify
the remaining assertion precisely: its right side
must equal the image of $C_M(A_5)$ for every $U$.
The finite permutation computation proves that
equality at the puncture-permutation level. It does
not establish it at every refinement. With only
exceptional orbits the discrete centralizer is
trivial, since the quotient sphere has just the
three differently ordered branch points and
$Z(A_5)=1$; the still unproved assertion is then
that the pure completed centralizer is trivial.

\emph{The minimal twelve-puncture case.}
There is a further exact localization of this remaining group.
Write $P=P\Gamma_{0,12}$ and choose $A=A_4\leq A_5$.
The action of $A$ on the twelve vertices is free and transitive.
Its quotient has one free puncture and three branch points of
orders $2,3,3$. The labelled commuting homeomorphism group on
the filled sphere is identified, by descent and the unique
commuting lift, with the identity homeomorphism component of
the quotient fixing those three branch points. It is
contractible. Evaluating the free quotient puncture consequently
identifies the mapping-class centralizer with the free group
$\pi_1(S^2\setminus\{3\text{ branch points}\})=F_2$.
Its permutation on the free fibre is the surjective monodromy
onto $A_4$. The pure centralizer is therefore the exact kernel
\[
 1\longrightarrow C_P(A)=F_{13}\longrightarrow F_2
                    \longrightarrow A_4\longrightarrow1.
                                                               \tag{F.7.8}
\]
Indeed it is the fundamental group of
$X=S^2\setminus R_A$, where the $A_4$ ramification set $R_A$
has $6+4+4=14$ points; its rank is also
$1+12(2-1)=13$ by Schreier. Label the vertices by $g\in A_4$
using a chosen vertex $z_0$. The inclusion is the configuration
map $z\mapsto(gz)_{g\in A_4}$, followed by the quotient by
M\"obius transformations. The proved solvable comparison and
full topology identify its completion with $C_{\widehat P}(A)$.
In particular the unknown group $C_{\widehat M}(A_5)$ lies in
this specified embedded $\widehat F_{13}$.

Each of the six subgroups $C_5\leq A_5$ has two fixed vertices
and two free five-cycles. Retain either one fixed vertex and
one five-cycle, giving twenty-four six-point sets, or retain
both fixed vertices and one five-cycle, giving twelve
seven-point sets. For each retained set $I$, the continuous
forgetting map $r_I$ is defined on $\widehat P$ and is equivariant
for the specified $C_5$. In either target the quotient has
exactly its two branch points and one free puncture. Its
discrete $C_5$ centralizer is therefore exactly the deck group
$C_5$: the quotient mapping class group fixing these three
decorated points is trivial. Prime-power comparison gives the
same completed centralizer. Its intersection with the pure
group is trivial, since the deck group permutes the five-cycle.
It follows that
\[
 C_{\widehat M}(A_5)\subset
 C_{\widehat P}(A)\cap
       \bigcap_{I\text{ among these }36\text{ sets}}\ker(r_I).
                                                               \tag{F.7.9}
\]

The six-point maps in (F.7.9) have a proved common kernel even
on the discrete $F_{13}$. The antipodal permutation on the
twelve vertices commutes with $A_5$. In the regular $A_4$
labelling it is right multiplication by an involution $h\in A_4$.
Every retained six-point set contains exactly one member of
each antipodal pair. Collisions of two entries $g_i z,g_j z$
at a ramification point with stabilizer $H\leq A_4$ occur
precisely when $g_j^{-1}g_i\in H$; these are the right cosets
$gH$. Thus there is no collision among the retained entries
at either fixed point of $h$. Each six-point configuration map
extends across both of these deleted points of $X$. Their
nontrivial peripheral meridians in $F_{13}$ map to the identity
under all twenty-four maps. Completion retains these meridians
as nontrivial primitive procyclic subgroups. Consequently the
joint six-point map on the embedded $\widehat F_{13}$ is
noninjective by these explicit elements.

Every seven-point set retains a collision at each of the
fourteen ramification points: seven choices from six pairs
force a pair, and seven choices from four triples force a
triple collision. The program verifies all twelve sets and
the preceding common invisible axis for the six-point sets,
using right cosets throughout. It also computes the induced
maps on $H_1(F_{13},\F_{101})$. A meridian at an $A_4$
ramification point maps to the sum of the pair generators
within its collision clusters. In the target pure sphere
group the pair generators satisfy the vertex-sum and total-sum
relations. The two fixed points of each cyclic stabilizer give
the same homology column. Each seven-point map and their joint
map have rank six, with seven-dimensional kernel; the joint
six-point map has rank five. These are exact calculations of
the indicated abelian quotients. They neither produce a
nontrivial $A_5$ centralizer element nor prove injectivity of
the joint seven-point map on $\widehat F_{13}$.

F.8 proves that all twelve seven-point kernels in (F.7.9)
are the same, both discretely and after completion. It further
proves that the discrete kernel contains no nonidentity reducible
mapping class and that the completed kernel meets every conjugate
of every full source cusp subgroup trivially. The exact four-point
characters detecting those cusps and all eighty-four finite
subset--inertia checks are supplied there. F.9 subsequently
constructs an explicit nontrivial common kernel element with
a full procyclic closure, and proves that all twenty-four
six-point restrictions factor through these seven-point maps.
The thirty-six-fold kernel is therefore nontrivial. The full
centralizer comparison and geometric conjugacy remain unproved;
no converse to the inclusion (F.7.9) is inferred.

\medskip\noindent\textbf{The exact common-$C_3$ comparison.}

The following maps give an exact model for the residual group in (F.7.1).
The fixed intersection displayed below is the remaining unresolved group.
All permutations use the five-letter identification of $A_5$ already
fixed in F.7--F.10.  Put
\[
 A=\operatorname{Stab}_{A_5}(5),\qquad
 r=(1\ 2\ 3),\qquad K=\langle r\rangle,
 \qquad\beta=(1\ 2)(4\ 5).
\]
Then $A\cong A_4$, $\beta K\beta^{-1}=K$ by inversion,
and $\langle A,\beta\rangle=A_5$.

\medskip\noindent\textbf{The actual quotient homomorphism.}
Let $P=P\Gamma_{0,12}$ at the icosahedral configuration.
The discrete branched-cover sequence for the $K$ action is
\[
 1\longrightarrow K\longrightarrow N(K)
       \longrightarrow Q\longrightarrow1,
\]
where $Q$ is the finite-index liftable subgroup of the
mapping class group of the quotient sphere with two decorated
branch points and four free punctures.  The centralizer
$C_P(K)$ has finite index in $N(K)$: both the centralizing
condition and the pure puncture condition have finite index.
Its intersection with $K$ is trivial, since nonidentity deck
maps permute the twelve punctures.  Its image $J$ in $Q$ is
therefore isomorphic to $C_P(K)$ and has finite index in $Q$.
Every element of $J$ fixes the four free quotient punctures.
It also fixes the two branch points separately: exchanging
them conjugates the order-three local deck generator to its
inverse, and hence cannot centralize it.  Thus
\[
             C_P(K)\xrightarrow{\ \sim\ }J
                         \leq P\Gamma_{0,6},
               \qquad [P\Gamma_{0,6}:J]<\infty.          \tag{F.7.10}
\]
The prime-power comparison and full induced topology proved
in F.7 identify $C_{\widehat P}(K)$ with
$\widehat{C_P(K)}$.  Completing the finite-index embedding
(F.7.10) identifies it with the specified open subgroup
$\widehat J\leq\widehat{P\Gamma}_{0,6}$.  This gives an
injective continuous homomorphism on the actual completed
centralizer; no completed intersection identity is used.

By F.7.8 the orbit map identifies $C_P(A)$ with the specified
$F_{13}=\pi_1(X)$, where
\[
 X=\mathbf P^1\setminus R_A,\qquad
 R_A=\{\infty\}\cup\{z:z(z^4-1)(z^8+14z^4+1)=0\}.
\]
Since $K\leq A$, composition with (F.7.10) is the injection
$j:F_{13}\hookrightarrow P\Gamma_{0,6}$ described next.
The solvable comparison and full topology for $A$ and $K$
show that its continuous extension
\[
                \widehat j:\widehat F_{13}
                       \hookrightarrow\widehat{P\Gamma}_{0,6}
                                                               \tag{F.7.11}
\]
is injective.  This is a different, explicitly specified map
from the thirty-six forgetful maps whose common kernel was
computed in F.9.

\medskip\noindent\textbf{The six labelled quotient points.}
In the F.9 tetrahedral coordinate,
\[
 r(z)=\frac{z+i}{z-i},\qquad
 q(z)=z+r(z)+r^2(z)
     =\frac{z^3+3iz+1-i}{(z-1)(z-i)}.
\]
The degree-three quotient $q$ is invariant under $r$.
Its branch points upstairs are
\[
 p_\pm=\frac{(1+i)(1\pm\sqrt3)}2,
\]
and their images are $3p_\pm$.  In the twelve-point
labelling of F.8 the four free $K$-orbits, in order, are
\[
 (0,8,1),\quad(11,3,5),\quad(2,6,9),\quad(7,10,4).
                                                               \tag{F.7.12}
\]
Their representatives have coordinates $z,1/z,-1/z,-z$.
Consequently the unnormalized labelled quotient map is
\[
 z\longmapsto
   (3p_+,3p_-,q(z),q(1/z),q(-1/z),q(-z)).                \tag{F.7.13}
\]
Its entries are distinct on $X$: any equality between two
free quotient points would identify their points in the
$A$ orbit, and a collision with a branch point would give a
nontrivial stabilizer in $A$.

Put $U(z)=(z-p_+)/(z-p_-)$.
Direct substitution gives
\[
       \frac{q(z)-3p_+}{q(z)-3p_-}=U(z)^3.
\]
Normalizing the first, second, and third points in (F.7.13)
to $0,\infty,1$ gives the actual map
\[
 \Sigma:X\longrightarrow M_{0,6},\qquad
 z\longmapsto(s_1(z),s_2(z),s_3(z)),
\]
where
\[
 s_1=\left(\frac{U(1/z)}{U(z)}\right)^3,
 \qquad
 s_2=\left(\frac{U(-1/z)}{U(z)}\right)^3,
 \qquad
 s_3=\left(\frac{U(-z)}{U(z)}\right)^3.                 \tag{F.7.14}
\]
The homomorphism $j$ in (F.7.11) is precisely $\Sigma_*$,
by descent of the same labelled orbit path.

\medskip\noindent\textbf{The second $A_4$ and the exact based involution.}
The involution $\beta$ exchanges the two branch points and
permutes the four free orbits in (F.7.12) by $(0\ 1)(2\ 3)$.
On the normalized moduli coordinates its relabelling is
therefore
\[
 B(s_1,s_2,s_3)=(s_1,s_1/s_3,s_1/s_2).                 \tag{F.7.15}
\]
For example, after swapping the branch points the new
normalizing coordinate is $u\mapsto s_1/u$, since the
first free point has become the old point $s_1$.
This derives (F.7.15) from the actual labelled quotient,
rather than selecting a conjugate linear action.

Use $z_*$ and the path
$\gamma:z_*\to5+z_*\to b=5+4i$ from (F.10.12).
The point $\Sigma(z_*)$ is fixed by $B$.  Let
$\tau=\Sigma\circ\gamma$, with the indicated direction.
For a chronological loop $\ell$ based at $\Sigma(b)$, the
transported involution is exactly
\[
 \beta_b(\ell)=
 \tau^{-1}*B(\tau)*B(\ell)*B(\tau)^{-1}*\tau.          \tag{F.7.16}
\]
The path endpoints in this expression match in order;
$B^2=1$ and $B\Sigma(z_*)=\Sigma(z_*)$ show directly that
$\beta_b^2=1$.  It induces the conjugation action of the
same transported finite mapping class $\beta$ used in F.10
on the injective quotient image (F.7.10).

Since $A$ and $\beta$ generate $A_5$, the exact residual
comparison is
\[
 \widehat j\bigl(C_{\widehat P}(A_5)\bigr)
   =\widehat j(\widehat F_{13})\cap
       \operatorname{Fix}\!\left(
          \beta_b:\widehat{P\Gamma}_{0,6}
                    \longrightarrow\widehat{P\Gamma}_{0,6}
                          \right).                    \tag{F.7.17}
\]
Indeed centralizing $A_5$ is exactly centralizing $A$ and
$\beta$.  For elements centralizing $A$, equation (F.7.11)
is injective; also $\beta$ normalizes $K$, so both the
element and its $\beta$ conjugate lie in the domain of
that injection.  Equality of their quotient images is
therefore equivalent to their equality upstairs.

\medskip\noindent\textbf{Executed exact checks and their scope.}
The program \path{code/part_a/second_a4_c3_attack.py}
uses exact arithmetic in $\Q(i,\sqrt5,\sqrt{\varphi+2})$
to check the $60$ M\"obius matrices, $180$ generator
products, and $720$ vertex-label images.  It verifies all
four orbits in (F.7.12), all representative functions,
the quotient identity, and (F.7.14).  It also checks
the fixed-locus calculation
\[
 \Sigma(z)\in\operatorname{Fix}(B)
 \quad\Longleftrightarrow\quad s_1(z)=s_2(z)s_3(z),
\]
whose numerator, up to a nonzero scalar and after
cancelling the displayed rational denominators, is
\[
 z(z^2+1)(z^8+16z^6+46z^4+16z^2+1).                    \tag{F.7.18}
\]
It verifies exactly that $z_*$ satisfies the relevant
quartic factor.  Finite intersection of these geometric
loci does not by itself calculate the intersection of
their completed fundamental groups in (F.7.17).

The auxiliary numerical braid discovery in
\path{work/part_a/f9_c3_johnson_probe.py} currently gives
no exclusion at the first nonzero outer Johnson degree.
It is explicitly exploratory: the crossings there have
not been certified for the actual based map (F.7.16),
and a higher-degree comparison must retain that path.
No conclusion about (F.7.17) uses that calculation.

This audit concerns the comparison requested for
\S30, folios 194--196, scans 383--387. F.6's
arithmetic open-image proof applies to $A_5$ without
using (F.7.1), so that proved arithmetic assertion
does not depend on the unresolved intersection.
\end{proof}
